\documentclass[11pt,letterpaper]{article}
\usepackage{physical_ai_legislation}
\usepackage[backend=biber,style=numeric,sorting=none]{biblatex}
\addbibresource{physical_ai_data_transparency.bib}

\title{%
  \textbf{Senate Bill 892}\\[6pt]
  \large California Physical AI Clinical Data Protection\\
  and Transparency Act of 2026\\[6pt]
  \normalsize An act to add Chapter 3.5 (commencing with Section 1798.300)\\
  to Title 1.81.5 of Part 4 of Division 3 of the Civil Code,\\
  relating to data protection in physical AI clinical trials.
}
\author{%
  Kevin Kawchak\\
  CEO, ChemicalQDevice\\
  \texttt{kevink@chemicalqdevice.com}
}
\date{March 24, 2026}

\begin{document}
\maketitle
\thispagestyle{fancy}

\begin{center}
\footnotesize
\textit{Unofficial independent draft derived from prior legislation;
no third-party endorsement, sponsorship, affiliation, or authorization
is expressed or implied; and was adapted using Claude Code Opus 4.6.}
\end{center}

\vspace{0.5em}

\begin{center}
\footnotesize
\textit{This work was adapted from original CFR documents in the public domain.
The original ICH document is copyrighted and may be used, reproduced,
incorporated into other works, adapted, modified, translated or distributed
under a public license. This current work is not endorsed or sponsored by
CFR, ICH, or FDA; and was adapted using Claude Code Opus 4.6.}
\end{center}

\vspace{1em}
\tableofcontents
\newpage

\section{Legislative Findings and Declarations}

The Legislature finds and declares all of the following:

\begin{enumerate}[label=(\alph*)]
\item Physical AI oncology clinical trials generate unprecedented volumes of
  patient data, including robotic telemetry, AI decision logs, digital twin
  state data, medical imaging, vital signs captured at minute-level resolution,
  and federated learning model parameters~\cite{newtrial2026}.

\item The 24-hour on-demand simulation generated 178 output files across 25
  sequential commits, including hourly patient records with minute-level vital
  sign data for 168 patients, detailed robot telemetry for 29 instances across
  24 hours, and PSL scoring data for all 10 robot
  types~\cite{newtrial2026}.

\item The MCP-PAI standard establishes a five-server interoperability topology
  (authorization, FHIR, DICOM, ledger, provenance) with 23 tools, six actor
  roles, five conformance levels, and hash-chained audit trails, providing a
  comprehensive data governance architecture~\cite{iche6r3adapt}.

\item Training data transparency requirements established by AB 2013 (2024)
  apply to the generative AI and VLA models integrated into Physical AI
  systems, and additional transparency requirements specific to robotic
  telemetry and clinical decision data are needed.

\item The cancer patient journey demonstrated federated learning across 12 sites
  using secure multi-party computation with differential privacy at epsilon
  equals 1.0, providing a model for privacy-preserving multi-site data
  sharing~\cite{patientjourney2026}.

\item Both specific details by the minute and details over 24 hours were
  captured and managed correctly by AI systems, creating a data governance
  challenge that requires clear rules for data retention, access, sharing, and
  deletion across multiple timescales.
\end{enumerate}

\section{Definitions}

For purposes of this chapter, the following definitions apply:

\begin{enumerate}[label=(\alph*)]
\item ``Physical AI clinical data'' means any data generated, collected,
  processed, or stored by Physical AI systems during oncology clinical trials,
  including robotic telemetry data, AI decision logs, digital twin state data,
  patient vital signs, medical imaging data, procedural records, adverse event
  records, and PSL and USL scoring data.

\item ``Robotic telemetry data'' means real-time operational data from Physical
  AI systems including position, force, velocity, temperature, power
  consumption, active or standby status, and performance metrics.

\item ``AI decision log'' means a record of decisions made by AI algorithms
  integrated into Physical AI systems, including treatment recommendations,
  scheduling decisions, resource allocation, and safety interventions.

\item ``Digital twin data'' means patient-specific computational model data
  including anatomical models, treatment simulation parameters, procedure
  rehearsal records, and predictive model outputs.

\item ``Hash-chained audit trail'' means an immutable, cryptographically linked
  sequence of data access and modification records using SHA-256 or equivalent
  algorithms, as specified in the adapted ICH E6(R3)~\cite{iche6r3adapt}.

\item ``Federated learning data'' means model parameters, gradient updates, and
  aggregated statistics shared between Physical AI systems across multiple
  clinical trial sites using privacy-preserving computation methods.

\item ``De-identification'' means the process of removing or obscuring personal
  health information using the HIPAA Safe Harbor methodology with 18
  identifier categories, as required by the adapted regulatory
  frameworks~\cite{cfr50adapt}.
\end{enumerate}

\section{Data Collection Transparency}

\subsection{Notice Requirements}

\begin{enumerate}[label=(\alph*)]
\item Before enrollment, each patient shall receive a comprehensive data
  collection notice specifying all categories of Physical AI clinical data that
  will be generated during their participation, the purposes for which each
  category will be used, and the retention period for each category.

\item The notice shall identify the five MCP-PAI servers that will process
  patient data (authorization, FHIR, DICOM, ledger, provenance) and describe
  the role of each server in plain language~\cite{iche6r3adapt}.

\item The notice shall disclose whether patient data will be used for federated
  learning and, if so, describe the differential privacy protections applied,
  including the privacy budget (epsilon value)~\cite{patientjourney2026}.

\item The notice shall be updated whenever new categories of data collection are
  introduced or existing categories are materially changed.
\end{enumerate}

\subsection{Real-Time Data Access}

\begin{enumerate}[label=(\alph*)]
\item Patients shall have the right to access their Physical AI clinical data
  in real time through a secure patient portal, including current vital signs,
  procedure status, robot telemetry relevant to their care, and digital twin
  visualizations.

\item The site shall provide data access in standardized formats, including
  FHIR for clinical data and DICOM for medical imaging data.

\item Data access requests shall be fulfilled within 24 hours for historical
  data not available through the real-time portal.
\end{enumerate}

\section{Data Protection Standards}

\subsection{Encryption and Access Control}

\begin{enumerate}[label=(\alph*)]
\item All Physical AI clinical data shall be encrypted at rest using AES-256 or
  equivalent and in transit using TLS 1.3 or later.

\item Access to Physical AI clinical data shall follow deny-by-default
  authorization policies, meaning no access is permitted unless explicitly
  authorized for a specific purpose, role, and time period~\cite{iche6r3adapt}.

\item The six MCP-PAI actor roles (patient, clinician, investigator, sponsor,
  regulator, system) shall define the maximum data access scope for each
  role, with no role having unrestricted access to all data
  categories~\cite{iche6r3adapt}.

\item Multi-factor authentication shall be required for all human access to
  Physical AI clinical data systems.
\end{enumerate}

\subsection{Audit Trail Requirements}

\begin{enumerate}[label=(\alph*)]
\item Hash-chained audit trails shall be maintained for all Physical AI clinical
  data transactions, providing cryptographic proof of data integrity and an
  immutable record of who accessed what data, when, and for what
  purpose~\cite{iche6r3adapt}.

\item Audit trail records shall be retained for a minimum of 15 years following
  the completion of the clinical trial, or longer if required by applicable
  regulations.

\item The site shall conduct quarterly audit trail integrity verification using
  cryptographic chain validation.

\item Audit trail data shall be stored in a separate system from clinical data
  to prevent simultaneous compromise.
\end{enumerate}

\subsection{De-identification Standards}

\begin{enumerate}[label=(\alph*)]
\item All Physical AI clinical data shared outside the treating site shall be
  de-identified using the HIPAA Safe Harbor methodology, removing all 18
  categories of identifiers~\cite{cfr50adapt}.

\item Robotic telemetry data shall be assessed for re-identification risk
  before sharing, as movement patterns and procedural signatures may
  constitute quasi-identifiers.

\item Digital twin data shall be fully de-identified before any use in
  federated learning, research, or quality improvement activities.

\item The site shall employ a qualified statistician or privacy officer to
  verify that de-identified data sets meet the threshold for non-identifiability.
\end{enumerate}

\section{Federated Learning and Multi-Site\\Data Sharing}

\begin{enumerate}[label=(\alph*)]
\item Physical AI oncology clinical trial sites that participate in federated
  learning shall implement the privacy-preserving computation methods
  demonstrated in the cancer patient journey, including secure multi-party
  computation with differential privacy~\cite{patientjourney2026}.

\item The privacy budget (epsilon value) for differential privacy shall not
  exceed 1.0 per federated learning round without explicit patient consent and
  IRB approval.

\item Model parameters shared through federated learning shall be verified to
  not contain memorized patient data through membership inference testing.

\item Federated learning participants shall maintain data harmonization
  standards using DICOM, FHIR, ICD-10 to SNOMED CT mapping, and LOINC
  coding~\cite{iche6r3adapt}.

\item The simulation demonstrated that 70 federated learning rounds across 12
  sites achieved clinically significant model improvements while maintaining
  patient privacy, providing evidence that multi-site Physical AI data sharing
  can be conducted safely~\cite{patientjourney2026}.
\end{enumerate}

\section{Data Retention and Deletion}

\begin{enumerate}[label=(\alph*)]
\item Physical AI clinical data shall be retained for the duration required by
  applicable federal and state regulations, including a minimum of 2 years
  following investigational drug approval or 2 years after investigation
  discontinuation, consistent with 21 CFR Part 312~\cite{cfr312adapt}.

\item Upon trial completion and satisfaction of regulatory retention
  requirements, patients shall have the right to request deletion of their
  digital twin data and patient-specific data from Physical AI system local
  storage.

\item Deletion requests shall be fulfilled within 30 days using secure data
  erasure methods that comply with HIPAA Safe Harbor requirements, and the
  site shall provide written confirmation of deletion~\cite{cfr312adapt}.

\item Aggregated, de-identified data may be retained indefinitely for research
  purposes, provided it meets the de-identification standards established in
  Section 4.3.

\item When a site's authorization is revoked or a trial is terminated, all
  patient-specific data shall be securely erased from Physical AI system local
  storage within 30 days of decommissioning, consistent with the IND
  withdrawal requirements of the adapted 21 CFR Part 312
  \S~312.38~\cite{cfr312adapt}.
\end{enumerate}

\section{Cybersecurity Incident Reporting}

\begin{enumerate}[label=(\alph*)]
\item Any cybersecurity incident affecting Physical AI clinical data shall be
  reported to the department, to affected patients, and to the FDA in
  accordance with the following timelines:
  \begin{enumerate}[label=(\arabic*)]
  \item Within 7 days if the incident involves potential patient harm or
    breach of protected health information for 500 or more individuals.
  \item Within 15 days for all other cybersecurity incidents affecting Physical
    AI systems~\cite{cfr312adapt}.
  \end{enumerate}

\item The site shall maintain a cybersecurity incident response plan that
  addresses Physical AI-specific threats including robot control system
  compromise, AI model poisoning, digital twin data manipulation, and
  MCP-PAI server infrastructure attacks.

\item Annual cybersecurity assessments shall include penetration testing of
  Physical AI system interfaces, MCP-PAI server infrastructure, and network
  segmentation between clinical and administrative systems.

\item The simulation demonstrated continuous operation across 24 hours with no
  cybersecurity incidents, including secure management of 10 gigabits per
  second internal network traffic supporting 29 robot instances and patient
  data systems~\cite{sitespec2026}.
\end{enumerate}

\section{AI Model Transparency}

\begin{enumerate}[label=(\alph*)]
\item In accordance with the principles of AB 2013 (2024), sites shall maintain
  and make available upon request documentation of the AI models integrated
  into Physical AI systems, including model architecture, training data
  sources, validation methodology, and performance benchmarks.

\item The five categories of AI recognized in the adapted ICH E6(R3) -
  generative AI (VLA models, diffusion policies), agentic AI (LLM-based
  assistants), reinforcement learning, self-supervised learning, and
  supervised learning - shall each have category-specific transparency
  requirements~\cite{iche6r3adapt}.

\item Pre-determined change control plans (PCCPs) for AI model updates shall be
  filed with the department and made available to patients upon request,
  specifying which model changes require protocol amendments versus those
  within the PCCP scope~\cite{cfr312adapt}.

\item AI model performance shall be reported in the site's annual report
  alongside PSL and USL scores, providing a comprehensive transparency
  record~\cite{cfr312adapt}.
\end{enumerate}

\section{Enforcement and Penalties}

\begin{enumerate}[label=(\alph*)]
\item Violations of data protection standards established under this chapter
  shall be subject to civil penalties of not less than two thousand five
  hundred dollars (\$2,500) and not more than seven thousand five hundred
  dollars (\$7,500) per violation per patient affected.

\item Intentional or reckless failure to maintain hash-chained audit trails or
  to report cybersecurity incidents within required timelines shall constitute
  an aggravated violation subject to penalties of up to twenty-five thousand
  dollars (\$25,000) per violation.

\item The Attorney General may bring an action to enforce this chapter on behalf
  of the people of the state.

\item Nothing in this chapter shall be construed to limit any private right of
  action available to patients under the California Consumer Privacy Act or
  the California Privacy Rights Act.
\end{enumerate}

\section{Effective Date}

This act shall take effect on January 1, 2027.

\printbibliography[title={References}]

\end{document}
